\documentclass{article}
\usepackage{amsmath, amssymb, amsthm}
\usepackage{hyperref}
\usepackage{geometry}
\geometry{margin=1in}

\title{Erd\H{o}s Problem \#302}
\date{Last updated: 2025-08-31}
\author{Source: erdosproblems.com}

\begin{document}
\maketitle

\noindent\textbf{Status:} OPEN \quad \textbf{No prize}

\noindent\textbf{Tags:} number theory, unit fractions

\medskip

\noindent\textbf{Problem Statement:}

Let $f(N)$ be the size of the largest $A\subseteq \{1,\ldots,N\}$ such that there are no solutions to\[\frac{1}{a}= \frac{1}{b}+\frac{1}{c}\]with distinct $a,b,c\in A$? Estimate $f(N)$. In particular, is $f(N)=(\tfrac{1}{2}+o(1))N$?




The colouring version of this is [303] , which was solved by Brown and R\"{o}dl \cite{BrRo91}. One can take either $A$ to be all odd integers in $[1,N]$ or all integers in $[N/2,N]$ to show $f(N)\geq (1/2+o(1))N$. Wouter van Doorn has proved (see this note ) that\[f(N) \leq (9/10+o(1))N.\]Stijn Cambie has observed that\[f(N)\geq (5/8+o(1))N,\]taking $A$ to be all odd integers $\leq N/4$ and all integers in $[N/2,N]$. Stijn Cambie has also observed that, if we allow $b=c$, then there is a solution to this equation when $\lvert A\rvert \geq (\tfrac{2}{3}+o(1))N$, since then there must exist some $n,2n\in A$. See also [301] and [327] .




References

[BrRo91] Brown, Tom C. and R\"{o}dl, Voijtech, Monochromatic solutions to equations with unit fractions . Bull. Austral. Math. Soc. (1991), 387-392.

\noindent\textbf{Related OEIS sequences:} A390395


\bigskip
\noindent\small{Source: \url{https://www.erdosproblems.com/302}}
\end{document}
